\documentclass[12pt]{ximera}
\title{Homework 1: Algebra Review}
\begin{document}
\begin{abstract}
Algebra review before the main content of the course: expanding and factoring
polynomials, the zero-product property, and solving a linear equation.
\end{abstract}
\maketitle

\section*{Sections R.1--R.4: Algebra Review}

Here are a few problems for algebra review, before we get into the meat of the
semester. If any of these give you trouble, work through the hints and solutions ---
and bring the ones that still give you trouble to class.

\begin{problem}
Expand the polynomial $(6x+5)(3x-2)$ and combine like terms.

\[ (6x+5)(3x-2) = \answer{18x^2+3x-10} \]

\begin{hint}
Multiply every term of the first factor by every term of the second (FOIL), then
collect the two $x$ terms.
\end{hint}

\begin{feedback}
\begin{align*}
(6x+5)(3x-2) &= 6x\cdot 3x + 6x\cdot(-2) + 5\cdot 3x + 5\cdot(-2)\\
             &= 18x^2 - 12x + 15x - 10\\
             &= 18x^2 + 3x - 10.
\end{align*}
\end{feedback}
\end{problem}

\begin{problem}
Factor the polynomial $16x^2+64$ by bringing the greatest common factor out front.

Write the answer in the form (GCF) times a polynomial:
\[ 16x^2 + 64 = \answer{16}\left(\answer{x^2+4}\right) \]

\begin{hint}
What is the largest number that divides both $16$ and $64$? Note that $x$ is
\emph{not} a common factor here --- the second term has no $x$ in it.
\end{hint}

\begin{feedback}
The greatest common factor of $16$ and $64$ is $16$ itself, since $64=16\cdot 4$, and
$64$ has no factor of $x$, so the GCF is $16$:
\[ 16x^2 + 64 = 16\left(x^2 + 4\right). \]
\end{feedback}
\end{problem}

\begin{problem}
Determine which of the following is true for all real numbers, and give a
counterexample (that is, values of $A$, $B$, and $C$ for which it fails) if it is
false.

\begin{prompt}
\textbf{(a)} $\dfrac{A+B}{C} = \dfrac{A}{C} + \dfrac{B}{C}$

\begin{multipleChoice}
\choice[correct]{True for all real numbers (where $C\neq 0$)}
\choice{False}
\end{multipleChoice}

\begin{feedback}
This one is true. Dividing a sum by $C$ divides each term by $C$, which is exactly
what adding two fractions with the common denominator $C$ undoes:
$\frac{A}{C}+\frac{B}{C}=\frac{A+B}{C}$.
\end{feedback}
\end{prompt}

\begin{prompt}
\textbf{(b)} $\dfrac{A}{B+C} = \dfrac{A}{B} + \dfrac{A}{C}$

\begin{multipleChoice}
\choice{True for all real numbers (where the denominators are nonzero)}
\choice[correct]{False}
\end{multipleChoice}
\end{prompt}

\begin{prompt}
\textbf{(b, continued)} Give a counterexample. Take $A=1$, $B=1$, and $C=1$. Then the
left-hand side equals $\answer{0.5}$ and the right-hand side equals $\answer{2}$.

\begin{feedback}
With $A=B=C=1$ the left-hand side is $\frac{1}{1+1}=\frac12$, while the right-hand
side is $\frac11+\frac11=2$. Since $\frac12\neq 2$, the statement is false. A sum in
the \emph{denominator} cannot be split apart the way a sum in the numerator can.
\end{feedback}
\end{prompt}
\end{problem}

\begin{problem}
What are the possible values of $x$ if $(x-4)(x+2)=0$?

Enter the smaller value first: $x = \answer{-2}$ or $x = \answer{4}$.

\begin{hint}
If a product of two numbers is zero, at least one of the factors must be zero.
Set each factor equal to $0$ and solve.
\end{hint}

\begin{feedback}
By the zero-product property, either $x-4=0$ or $x+2=0$, giving $x=4$ or $x=-2$.
\end{feedback}
\end{problem}

\begin{problem}
Isolate $x$ to solve the equation
\[ \frac{x}{3} - 5 = 6. \]

$x = \answer{33}$

\begin{hint}
Undo the operations in the reverse order they were applied: first add $5$ to both
sides, then multiply both sides by $3$.
\end{hint}

\begin{feedback}
Adding $5$ to both sides gives $\frac{x}{3}=11$, and multiplying both sides by $3$
gives $x=33$. Check: $\frac{33}{3}-5=11-5=6$. \checkmark
\end{feedback}
\end{problem}


\end{document}
